\chapter{实体分类体系简表}\label{ux5b9eux4f53ux5206ux7c7bux4f53ux7cfbux7b80ux8868}

本研究在实体抽取阶段采用``双重标签''体系：第一层为面向地方文化记忆识别的 AI 分级标签，分为 6 个一级体系、28 个二级类型；第二层为便于同官方资源名录衔接的官方分类标签。前者用于保证典籍文本中的人物、空间、技艺、文献和历史时序能够被完整识别，后者用于和非遗名录、文物保护单位、历史文化名镇名村、公共文化资源等管理口径发生对应。

{\def\LTcaptype{none} % do not increment counter
\begin{longtable}{@{}
  >{\centering\arraybackslash}p{(\linewidth - 6\tabcolsep) * \real{0.2500}}
  >{\centering\arraybackslash}p{(\linewidth - 6\tabcolsep) * \real{0.2500}}
  >{\centering\arraybackslash}p{(\linewidth - 6\tabcolsep) * \real{0.2500}}
  >{\centering\arraybackslash}p{(\linewidth - 6\tabcolsep) * \real{0.2500}}@{}}
\toprule\noalign{}
\endhead
\bottomrule\noalign{}
\endlastfoot
\textbf{一级体系} & \textbf{二级类型} & \textbf{主要识别对象示例} & \textbf{可对应的官方分类口径} \\
A 非遗文化体系 & A1 表演艺术类非遗 & 醒狮、粤剧、龙舟说唱、十番音乐 & 非物质文化遗产：传统舞蹈、传统戏剧、曲艺、传统音乐等 \\
A 非遗文化体系 & A2 传统技艺类非遗 & 灰塑、藤编、木雕、砖雕、陶塑 & 非物质文化遗产：传统技艺、传统美术 \\
A 非遗文化体系 & A3 民俗节庆类非遗 & 生菜会、龙舟竞渡、庙会、花灯会 & 非物质文化遗产：民俗 \\
A 非遗文化体系 & A4 信俗礼仪类非遗 & 民间信仰、龙母诞、地方礼俗 & 非物质文化遗产：民俗 \\
A 非遗文化体系 & A5 传统体育游艺类非遗 & 洪拳、咏春、龙舟竞技、狮艺 & 非物质文化遗产：传统体育、游艺与杂技 \\
A 非遗文化体系 & A6 饮食酿造类非遗及文化物产 & 九江双蒸酒、西樵大饼 & 非物质文化遗产：传统技艺 \\
B 物质文化遗产体系 & B1 古建筑类 & 书院、祠堂、庙宇、楼阁、塔桥 & 不可移动文物：古建筑 \\
B 物质文化遗产体系 & B2 宗教建筑类 & 宝峰寺、南海观音寺、云泉仙馆 & 不可移动文物：古建筑 \\
B 物质文化遗产体系 & B3 纪念性建筑与名人故居类 & 康有为故居、黄飞鸿纪念馆 & 不可移动文物：近现代重要史迹及代表性建筑 \\
B 物质文化遗产体系 & B4 古遗址与生产遗存类 & 古窑址、聚落遗存 & 不可移动文物：古文化遗址 \\
B 物质文化遗产体系 & B5 石刻碑记类 & 碑刻、摩崖石刻、匾额题记 & 不可移动文物：石窟寺及石刻；可移动文物 \\
B 物质文化遗产体系 & B6 古村落与聚落遗产类 & 松塘村、烟桥村、仙岗村 & 历史文化名城名镇名村：历史文化名村、传统村落、历史建筑等 \\
C 传承主体体系 & C1 历史文化人物 & 康有为、湛若水、黄飞鸿 & 无直接官方管理分类，作为文化记忆实体保留 \\
C 传承主体体系 & C2 非遗传承人及技艺人物 & 工匠、艺师、传承人 & 非遗传承主体：传承人 \\
C 传承主体体系 & C3 文物营建与守护人物 & 创建者、修建者、捐建者、守护者 & 无直接官方管理分类，作为营建与保护关系实体保留 \\
C 传承主体体系 & C4 宗族姓氏与地方社群 & 九江关氏、松塘区氏、地方社群 & 非遗传承主体：传承群体 \\
D 文化空间体系 & D1 山川水系空间 & 西樵山、西江、北江 & 自然保护地：自然公园等 \\
D 文化空间体系 & D2 镇街圩市空间 & 九江镇、丹灶镇、西樵镇、传统圩市 & 历史文化名城名镇名村：历史文化名镇等 \\
D 文化空间体系 & D3 历史街区与传统片区 & 历史街区、传统片区、老街区 & 历史文化名城名镇名村：历史文化街区、历史建筑 \\
D 文化空间体系 & D4 传承场所与活动场地 & 传习所、武馆、训练场、活动场地 & 公共文化资源：公共文化设施、公共文化活动 \\
E 文献记忆体系 & E1 地方志类 & 《南海县志》《大德南海志》等 & 可移动文物：古代文物、近代现代文物 \\
E 文献记忆体系 & E2 族谱家乘类 & 族谱、家乘 & 可移动文物：古代文物、近代现代文物 \\
E 文献记忆体系 & E3 碑记题咏类 & 碑记、诗文题刻 & 不可移动文物；可移动文物 \\
E 文献记忆体系 & E4 文集著述类 & 文集、学术著作、地方著述 & 可移动文物：古代文物、近代现代文物 \\
E 文献记忆体系 & E5 口述史与地方记忆材料 & 口述史、传说、地方记忆材料 & 非物质文化遗产：民间文学 \\
F 历史时序体系 & F1 朝代年号类 & 明代、清代、民国、光绪二十四年 & 无直接官方管理分类，作为时序标注实体保留 \\
F 历史时序体系 & F2 历史事件类 & 公车上书、戊戌变法、地方历史事件 & 无直接官方管理分类，作为事件实体保留 \\
F 历史时序体系 & F3 发展阶段类 & 形成期、兴盛期、转型期等阶段表述 & 无直接官方管理分类，作为阶段判断实体保留 \\
\end{longtable}
}

\chapter{165 条文化载体样本分类与镇街分布统计}\label{ux6761ux6587ux5316ux8f7dux4f53ux6837ux672cux5206ux7c7bux4e0eux9547ux8857ux5206ux5e03ux7edfux8ba1}

按类型：不可移动文物保护单位 / 空间锚点 80、非遗空间锚点 36、文化景观 19、圩市街区 18、历史文化名村/传统村落 12。

按镇街：西樵镇 89、九江镇 33、丹灶镇 30、桂城街道 8、大沥镇 3、狮山镇 1、里水镇 1。当前指数表已统一到镇街字段，因此正文和附录均按同一镇街口径统计。

完整清单作为载体级指数附表留存，正文仅摘录类型分布、镇街分布和典型样本。

官方资源扩展候选库总量为 710 条，其中保留原锚点 220 条，新增去重后的官方候选锚点 490 条。该库用于附加诊断和后续人工校核，不直接替代本论文 165 条指数计算样本。

\chapter{载体级指数与错位分类样例}\label{ux8f7dux4f53ux7ea7ux6307ux6570ux4e0eux9519ux4f4dux5206ux7c7bux6837ux4f8b}

\textbf{C.1 沉睡潜力区（MI 最负 Top 10）}

{\def\LTcaptype{none} % do not increment counter
\begin{longtable}{@{}lccccc@{}}
\toprule\noalign{}
名称 & 镇街 & CMI & OAI & THI & MI \\
\midrule\noalign{}
\endhead
\bottomrule\noalign{}
\endlastfoot
西樵镇松塘村 & 西樵镇 & 80.39 & 100.0 & 4.48 & −85.72 \\
九江镇烟桥烟南村 & 九江镇 & 84.51 & 100.0 & 7.09 & −85.16 \\
丹灶镇仙岗社区 & 丹灶镇 & 69.35 & 100.0 & 4.48 & −80.20 \\
赛龙舟（九江传统龙舟） & 九江镇 & 78.11 & 75.0 & 0.00 & −76.55 \\
西樵山仙姑亭 & 西樵镇 & 98.55 & 50.0 & 0.00 & −74.28 \\
丹灶镇南沙棋盘村 & 丹灶镇 & 70.98 & 75.0 & 0.00 & −72.99 \\
丹灶镇棋盘村 & 丹灶镇 & 70.98 & 75.0 & 0.00 & −72.99 \\
九江镇烟桥村 & 九江镇 & 84.77 & 75.0 & 7.09 & −72.79 \\
西樵镇百西村头村 & 西樵镇 & 80.91 & 75.0 & 7.09 & −70.86 \\
丹灶镇仙岗村 & 丹灶镇 & 70.47 & 75.0 & 4.48 & −68.26 \\
\end{longtable}
}

\textbf{C.2 空心景点区（MI 最正 Top 10）}

{\def\LTcaptype{none} % do not increment counter
\begin{longtable}{@{}lccccc@{}}
\toprule\noalign{}
名称 & 镇街 & CMI & OAI & THI & MI \\
\midrule\noalign{}
\endhead
\bottomrule\noalign{}
\endlastfoot
南海农谚 & 桂城街道 & 0.00 & 25.0 & 87.87 & 75.37 \\
南海广式旺阁酱油酿造技艺 & 桂城街道 & 0.00 & 25.0 & 70.00 & 57.50 \\
大伸市 & 九江镇 & 0.00 & 0.0 & 54.82 & 54.82 \\
光分亭 & 西樵镇 & 0.00 & 50.0 & 77.57 & 52.57 \\
红棉公园民国建筑群 & 九江镇 & 0.00 & 0.0 & 51.30 & 51.30 \\
沙基圩 & 西樵镇 & 0.00 & 0.0 & 46.88 & 46.88 \\
湖山胜迹门楼 & 西樵镇 & 14.70 & 50.0 & 78.27 & 45.92 \\
禄舟万寿堂 & 西樵镇 & 0.00 & 50.0 & 67.96 & 42.96 \\
大桐圩 & 西樵镇 & 0.00 & 0.0 & 42.58 & 42.58 \\
西联村神诞 & 桂城街道 & 0.00 & 25.0 & 55.07 & 42.57 \\
\end{longtable}
}

\textbf{C.3 核心耦合区（全部样本）}

{\def\LTcaptype{none} % do not increment counter
\begin{longtable}{@{}
  >{\raggedright\arraybackslash}p{(\linewidth - 10\tabcolsep) * \real{0.1304}}
  >{\centering\arraybackslash}p{(\linewidth - 10\tabcolsep) * \real{0.1739}}
  >{\centering\arraybackslash}p{(\linewidth - 10\tabcolsep) * \real{0.1739}}
  >{\centering\arraybackslash}p{(\linewidth - 10\tabcolsep) * \real{0.1739}}
  >{\centering\arraybackslash}p{(\linewidth - 10\tabcolsep) * \real{0.1739}}
  >{\centering\arraybackslash}p{(\linewidth - 10\tabcolsep) * \real{0.1739}}@{}}
\toprule\noalign{}
\begin{minipage}[b]{\linewidth}\raggedright
名称
\end{minipage} & \begin{minipage}[b]{\linewidth}\centering
镇街
\end{minipage} & \begin{minipage}[b]{\linewidth}\centering
CMI
\end{minipage} & \begin{minipage}[b]{\linewidth}\centering
OAI
\end{minipage} & \begin{minipage}[b]{\linewidth}\centering
THI
\end{minipage} & \begin{minipage}[b]{\linewidth}\centering
MI
\end{minipage} \\
\midrule\noalign{}
\endhead
\bottomrule\noalign{}
\endlastfoot
云泉仙馆 & 西樵镇 & 70.25 & 75.0 & 77.78 & 5.15 \\
西樵山百步云梯 & 西樵镇 & 98.58 & 50.0 & 66.51 & −7.78 \\
九江吴家大院 & 九江镇 & 69.28 & 75.0 & 66.04 & −6.10 \\
西樵山抗日阵亡将士暨死难同胞纪念碑 & 西樵镇 & 98.61 & 50.0 & 65.58 & −8.73 \\
平洲传统玉器制作技艺 & 桂城街道 & 57.28 & 50.0 & 56.63 & 2.99 \\
\end{longtable}
}

\chapter{潜力释放条件相关矩阵}\label{ux6f5cux529bux91caux653eux6761ux4ef6ux76f8ux5173ux77e9ux9635}

\textbf{D.1 载体级 Pearson 矩阵（n = 165）} 见第 6.3.1 节表 6.4。

\textbf{D.2 镇街级 Pearson 矩阵（n = 7）} 见第 6.3.2 节表 6.5。

\textbf{D.3 载体级 Spearman 矩阵与镇街级 Spearman 矩阵} 见补充附表。

\textbf{D.4 关键双变量对比}

\begin{itemize}
\item
  物质载体数 × POI 总量（镇街级）: Pearson r = −0.475；
\item
  非遗项目数 × POI 总量（镇街级）: Pearson r = 0.372；
\item
  CMI × MI（载体级）: Pearson r = −0.495（正文核心发现）；
\item
  OAI × MI（载体级）: Pearson r = −0.440（正文核心发现）；
\item
  OAI\_mean × THI\_mean（镇街级）: Pearson r = 0.702（评级与 THI 的关联）。
\end{itemize}
